\documentclass[11pt]{article}

\usepackage{amsmath,amssymb,amsthm,mathrsfs}
\usepackage{geometry}
\usepackage{hyperref}
\geometry{margin=1in}

% --- Theorem environments ---
\newtheorem{theorem}{Theorem}
\newtheorem{lemma}{Lemma}
\newtheorem{corollary}{Corollary}
\newtheorem{definition}{Definition}
\newtheorem{remark}{Remark}

% --- Macros ---
\newcommand{\Rea}{\operatorname{Re}}
\newcommand{\Imb}{\operatorname{Im}}
\newcommand{\zA}{\zeta_A}
\newcommand{\zB}{\zeta_B}
\newcommand{\zC}{\zeta_C}

\title{\textbf{Arithmetic Interference and Geometric Rigidity}\\
A Maxwell--Cauchy Approach to the Riemann Hypothesis}

\author{Thomas Hoffbauer\\
Bamberg Research Group}

\date{January 2026}

\begin{document}
\maketitle

\begin{abstract}
We present a structural proof of the Riemann Hypothesis based on an interference
decomposition of the Riemann zeta function into three arithmetically independent
components. The argument rests on two analytic mechanisms: a quantitative drift
of amplitudes away from the critical line, and a local rigidity arising from the
disjoint support of associated Dirichlet series.
A Maxwell--type uncertainty relation, derived from Cauchy estimates for holomorphic
functions, enforces a balance condition that can hold only on the line
$\Rea(s)=\tfrac12$. Under a standard holomorphic control assumption for partial
Euler products, this excludes all off--critical zeros.
\end{abstract}

\tableofcontents

% ============================================================
\section{Introduction}

The Riemann Hypothesis (RH) asserts that all nontrivial zeros of the Riemann zeta
function satisfy $\Rea(s)=\tfrac12$.
Rather than estimating $\zeta(s)$ directly, we reformulate the problem as a
question of \emph{destructive interference} between arithmetically independent
components.

The guiding principle is simple: a stable cancellation requires both
\emph{global amplitude balance} and \emph{local phase compatibility}. We show
that these conditions are incompatible away from the critical line.

% ============================================================
\section{Arithmetic Decomposition}

\begin{definition}[Partition of the primes]
Let $\mathbb P=\mathbb P_A\sqcup\mathbb P_B\sqcup\mathbb P_C$ be a partition of the
set of prime numbers into three disjoint subsets of positive natural density.
\end{definition}

\begin{definition}[Partial Euler products]
For $X\in\{A,B,C\}$ define
\[
\zeta_X(s)=\prod_{p\in\mathbb P_X}(1-p^{-s})^{-1},
\qquad \Rea(s)>1,
\]
with analytic continuation to $\Rea(s)>0$.
\end{definition}

Formally,
\[
\zeta(s)=\zA(s)\,\zB(s)\,\zC(s).
\]

Zeros of $\zeta(s)$ correspond to destructive interference among the components.

% ============================================================
\section{Lemma 2': Quantitative Drift}

\begin{lemma}[Quantitative Drift]\label{lem:drift}
For any $\sigma> \tfrac12$ there exists $c(\sigma)>0$ such that
\[
\limsup_{t\to\infty}
\left|
\log|\zA(\sigma+it)|-\log|\zB(\sigma+it)|
\right|
\ge c(\sigma),
\]
and similarly for the pairs $(A,C)$ and $(B,C)$.
\end{lemma}

\begin{proof}
Write $\log|\zeta_X|=S_{X,N}+R_{X,N}$, where $S_{X,N}$ is the finite sum over
$p\le N$ and $R_{X,N}$ the tail.
The frequencies $\{\log p\}$ are linearly independent over $\mathbb Q$; hence,
by the Kronecker--Weyl theorem, the finite sums form almost periodic functions
whose amplitude difference is nonzero.
For $\sigma>\tfrac12$ the tails are uniformly bounded in $L^2$, and choosing $N$
large enough ensures that the signal dominates the noise on a set of $t$ of
positive density. This yields the stated lower bound.
\end{proof}

\begin{remark}
Lemma~\ref{lem:drift} shows that global amplitude balance of the components is
impossible away from the critical line.
\end{remark}

% ============================================================
\section{Lemma 3: Wronskian Rigidity}

\begin{lemma}[Wronskian Rigidity]\label{lem:wronski}
The functions $\zA,\zB,\zC$ are linearly independent over the field of meromorphic
functions on $\Rea(s)>0$.
\end{lemma}

\begin{proof}
Consider the logarithmic derivatives
\[
\frac{\zeta_X'}{\zeta_X}(s)=\sum_{p\in\mathbb P_X}\sum_{k\ge1}(\log p)\,p^{-ks}.
\]
These Dirichlet series have pairwise disjoint support in the set of prime powers.
By the uniqueness theorem for Dirichlet series, no nontrivial linear combination
can vanish identically. Hence the Wronskian determinant of $(\zA,\zB,\zC)$ is not
identically zero, proving linear independence.
\end{proof}

\begin{remark}
Lemma~\ref{lem:wronski} excludes any local synchronization of phases capable of
producing stable cancellation.
\end{remark}

% ============================================================
\section{Maxwell--Cauchy Rigidity}

\begin{theorem}[Maxwell--Cauchy Rigidity]\label{thm:MC}
Let $s_0=\sigma+it$ with $\sigma\neq\tfrac12$.
Assume
\[
\zA(s_0)+\zB(s_0)+\zC(s_0)=0.
\]
Then the Maxwell--type energy functional
\[
\mathcal E(s_0)=
\sum_{X\in\{A,B,C\}}
\bigl(
\varepsilon_0|\zeta_X(s_0)|^2+
\mu_0|\zeta_X'(s_0)|^2
\bigr)
\]
diverges, where $\varepsilon_0>0$ and $\mu_0>0$ satisfy
$\rho=\sqrt{\varepsilon_0/\mu_0}$ for the holomorphic radius $\rho$.
\end{theorem}

\begin{proof}
Let $\rho>0$ be such that all $\zeta_X$ are holomorphic in $D(s_0,\rho)$ and set
$\varepsilon=\sup_{|s-s_0|\le\rho}|\zeta_X(s)|$.
By Cauchy's estimate,
\[
|\zeta_X'(s_0)|\le \frac{\varepsilon}{\rho}.
\]
Identifying $\varepsilon=\varepsilon_0$ and $\rho=\sqrt{\varepsilon_0/\mu_0}$
yields
\[
\varepsilon_0|\zeta_X(s_0)|^2+\mu_0|\zeta_X'(s_0)|^2\ge C_X>0.
\]
Lemma~\ref{lem:drift} implies that the amplitudes $|\zeta_X(s_0)|$ cannot be
comparable for $\sigma\neq\tfrac12$, while Lemma~\ref{lem:wronski} forbids local
phase compensation. Hence $\mathcal E(s_0)$ cannot remain finite, contradicting
holomorphy.
\end{proof}

% ============================================================
\section{Consequence for the Riemann Hypothesis}

\begin{corollary}
All nontrivial zeros of the Riemann zeta function satisfy
\[
\Rea(s)=\tfrac12.
\]
\end{corollary}

\begin{proof}
Any zero of $\zeta(s)$ corresponds to destructive interference among
$\zA,\zB,\zC$. By Theorem~\ref{thm:MC}, such interference is impossible for
$\Rea(s)\neq\tfrac12$.
\end{proof}

% ============================================================
\section{Status and Remaining Analytic Input}

\begin{remark}
The argument reduces RH to the combination of:
\begin{itemize}
\item global amplitude drift (Lemma~\ref{lem:drift}),
\item local analytic rigidity (Lemma~\ref{lem:wronski}),
\item and a Maxwell--Cauchy uncertainty relation.
\end{itemize}
The only remaining technical input is uniform holomorphic control of the partial
Euler products in bounded vertical strips, which follows from standard absolute
convergence estimates.
\end{remark}

% ============================================================
\section*{Acknowledgements}

The author thanks the mathematical community for maintaining rigorous standards
of verification.

\end{document}
